% Compile from teaching/ with the supplied customIMTbeamer directory.
% The user's attached class, themes and logos are retained unchanged.
\documentclass[table]{customIMTbeamer/imt_pres}
\usepackage{amsmath,amssymb,booktabs,listings}
\usetikzlibrary{arrows.meta,positioning}
\lstset{basicstyle=\ttfamily\small,breaklines=true,columns=fullflexible,showstringspaces=false}
\title[STTRPSD]{Truck and trailer routing\\with satellite depots}
\subtitle{STTRPSD --- Metaheuristics course project}
\author{Andon Tchechmedjiev}
\institute{IMT Mines Al\`es}
\date{Three sessions}
\hypersetup{pdftitle={STTRPSD metaheuristics course project},
 pdfauthor={Andon Tchechmedjiev},pdfsubject={Combinatorial optimization}}
\begin{document}
\imttitlepage

\section*{The routing problem}
\begin{frame}{One vehicle and two kinds of trip}
\begin{columns}[T,onlytextwidth]
\begin{column}{0.56\textwidth}
\centering
\begin{tikzpicture}[x=0.54cm,y=0.65cm,>=Stealth,
 dep/.style={rectangle,draw=imtDarkBlue,fill=imtDarkBlue,text=white,minimum size=6mm,font=\small},
 sat/.style={rectangle,draw=imtBlue,fill=imtGrey,minimum size=6mm,font=\small},
 cust/.style={circle,draw=imtDarkBlue,fill=white,minimum size=6mm,font=\small}]
\node[dep] (d) at (0,0) {0};
\node[sat] (s1) at (2,0) {1};
\node[sat] (s2) at (8,0) {2};
\node[cust] (c3) at (2,3) {3};
\node[cust] (c4) at (4.5,3) {4};
\node[cust] (c5) at (8,4) {5};
\node[cust] (c6) at (11,4) {6};
\draw[imtDarkBlue,very thick,->] (d)--(s1);
\draw[imtDarkBlue,very thick,->] (s1)--(s2);
\draw[imtDarkBlue,very thick,->] (s2.south) to[bend left=28] (d.south);
\draw[imtBlue,thick,->] (s1)--(c3);
\draw[imtBlue,thick,->] (c3)--(c4);
\draw[imtBlue,thick,->] (c4)--(s1);
\draw[imtBlue,thick,->] (s2) to[bend left=9] (c5);
\draw[imtBlue,thick,->] (c5) to[bend left=9] (s2);
\draw[imtRed,thick,->] (s2) to[bend right=8] (c6);
\draw[imtRed,thick,->] (c6) to[bend right=8] (s2);
\node[font=\scriptsize,anchor=south] at (2,3.35) {$q_3=6$};
\node[font=\scriptsize,anchor=south] at (4.5,3.35) {$q_4=4$};
\node[font=\scriptsize,anchor=south] at (8,4.35) {$q_5=8$};
\node[font=\scriptsize,anchor=south] at (11,4.35) {$q_6=7$};
\node[font=\scriptsize,anchor=south] at (0,0.6) {Main depot};
\node[font=\scriptsize,anchor=north] at (5,-1.35) {Dark: truck + trailer\quad Colour: truck only};
\end{tikzpicture}
\par\small Truck capacity $Q_V=10$: the trips carry $10$, $8$, and $7$.
\end{column}
\begin{column}{0.41\textwidth}
\begin{block}{The vehicle's journey}
\begin{itemize}
\item Truck and trailer visit selected satellites.
\item The trailer stays parked while the truck serves customers.
\item Each truck trip returns to its starting satellite.
\item Every customer is served exactly once.
\end{itemize}
\end{block}
\end{column}
\end{columns}
\end{frame}

\begin{frame}{The optimization problem}
\begin{block}{Minimize the complete distance}
\[
\min\quad \underbrace{\text{main-tour distance}}_{\text{truck + trailer}}
+\underbrace{\text{sum of truck-trip distances}}_{\text{truck only}}
\]
\end{block}
Choose the satellites and their order, assign customers to trips, and order the customers within each trip.
\begin{itemize}
\item Each truck trip carries at most $Q_V$.
\item The main tour starts and ends at depot 0, visiting each chosen satellite once.
\item Distances are Euclidean, \textbf{without integer rounding}.
\item The scorer checks feasibility and recomputes the full distance.
\end{itemize}
{\tiny Villegas et al. (2010), \emph{Engineering Applications of Artificial Intelligence} 23, 780--794.\quad
\url{https://doi.org/10.1016/j.engappai.2010.01.013}\par}
\end{frame}

\section*{Project work}
\begin{frame}{Greedy construction, local search, then a metaheuristic}
\begin{columns}[T,onlytextwidth]
\begin{column}{0.315\textwidth}
\begin{block}{1. Greedy construction}
Implement a rule that constructs feasible routes.

Validate the constructor on the ten practice instances and record a baseline.
\end{block}
\end{column}
\begin{column}{0.315\textwidth}
\begin{block}{2. Local search}
Improve routes with relocation, swap, 2-opt, or other justified moves.

Compare the results with the greedy constructor.
\end{block}
\end{column}
\begin{column}{0.315\textwidth}
\begin{block}{3. Metaheuristic}
Implement a genetic, population-based, swarm, or other justified metaheuristic.

Compare all three stages under the \textbf{30-second budget}.
\end{block}
\end{column}
\end{columns}
Reuse your constructor and local-search operators. The paper you pitched can guide the final method.
\begin{alertblock}{Project focus}
Design and evaluate a metaheuristic. The statement's integer model is an optional aid to understanding the constraints.
\end{alertblock}
\end{frame}

\begin{frame}[fragile]{The solver interface and benchmark}
\begin{columns}[T,onlytextwidth]
\begin{column}{0.49\textwidth}
\begin{block}{A common interface}
\begin{lstlisting}
solve --instance instance.json
      --solution solution.json
      --seed 0 --time-limit 30
\end{lstlisting}
Solvers can use any programming language.
\end{block}
\begin{exampleblock}{A main tour and rooted truck trips}
\[
\begin{aligned}
M&=(0,1,2,0),\\
r_1&=(1;3,4),\\
r_2&=(2;5),\qquad r_3=(2;6).
\end{aligned}
\]
Trip format: (satellite; customers).
\end{exampleblock}
\end{column}
\begin{column}{0.48\textwidth}
\begin{itemize}
\item \textbf{10 practice instances} for every team.
\item \textbf{10 other assessment instances} with similar characteristics.
\item \textbf{30 seconds per run}, including startup.
\item Three common seeds and one CPU core.
\item Failed or invalid runs earn zero for that run.
\end{itemize}
\end{column}
\end{columns}
\end{frame}

\section*{Assessment}
\begin{frame}{Five-minute exchanges with every group}
\begin{columns}[T,onlytextwidth]
\begin{column}{0.62\textwidth}
\begin{imttable}{Q[l]Q[c]Q[c]Q[c]}
Exchange & S1 & S2 & S3\\
11:30--11:35 & G4 & G5 & G6\\
11:36--11:41 & G1 & G3 & G7\\
11:42--11:47 & G2 & G7 & G5\\
11:48--11:53 & G7 & G6 & G2\\
11:54--11:59 & G6 & G4 & G3\\
12:00--12:05 & G3 & G1 & G4\\
12:06--12:11 & G5 & G2 & G1\\
\end{imttable}
\end{column}
\begin{column}{0.35\textwidth}
\begin{alertblock}{During the exchange}
Every member answers.

No computers, phones, or notes.
\end{alertblock}
One minute between groups. Other groups continue working.
\end{column}
\end{columns}
Including the paper pitch, every group has mean position 4. No group repeats an oral position.
\end{frame}

\begin{frame}{The final mark out of 20}
\begin{columns}[T,onlytextwidth]
\begin{column}{0.62\textwidth}
\begin{imttable}{X[l]Q[c]Q[c]}
Component & Raw & Final\\
Participation $P$ & /20 & /10\\
Code $C$ & /4 & /2\\
Report $R$ & /8 & /4\\
Benchmark performance $K$ & /8 & /4\\
\end{imttable}
\end{column}
\begin{column}{0.35\textwidth}
\begin{block}{Combine the raw marks}
\[
G=\frac{P+C+R+K}{2}
\]
\end{block}
\begin{exampleblock}{Participation}
Paper pitch and three exchanges: \textbf{four marks out of~5}.
\end{exampleblock}
\end{column}
\end{columns}
\end{frame}

\begin{frame}{The report has five sections}
\begin{enumerate}
\item \textbf{Method description}: representation, construction, search, and stopping rule.
\item \textbf{Benchmarking and results}: protocol, repeated runs, tables, and figures.
\item \textbf{Analysis and conclusions}: explain the results, tradeoffs, and limitations.
\item \textbf{Instructions to run the programme}: dependencies and exact commands.
\item \textbf{One page on generative AI}: your interactions, checks, and each student's contribution.
\end{enumerate}
\end{frame}

\section*{Getting started}
\begin{frame}{Session 1: a greedy solver}
\begin{enumerate}
\item Score the supplied toy solution and account for every cost term.
\item Choose a greedy decision rule and implement your constructor.
\item Validate all ten practice instances and record the baseline results.
\item Explain each decision and identify a neighbourhood for Session 2.
\end{enumerate}
\begin{block}{Be ready to explain your method}
Every member should be able to explain the greedy rule and why the resulting routes are feasible.
\end{block}
\end{frame}
\end{document}
